% ============================================================================
% COURS 09 — EXERCICES MACHINE À COURANT CONTINU
% ============================================================================

\section{Exercices du chapitre}\label{sec:mcc-exercices}

\begin{TLExercise}{Exercice 1 — Point nominal d'une MCC}
Une machine à courant continu à aimants permanents fonctionne sous
$U=48\,\mathrm V$. Sa résistance d'induit vaut $R=0{,}40\,\Omega$, le courant
$I=8{,}0\,\mathrm A$ et la vitesse $\Omega=180\,\mathrm{rad\,s^{-1}}$.
\begin{enumerate}
  \item Calculer la force électromotrice $E$.
  \item En déduire la constante $K_e$.
  \item Donner la valeur numérique de $K_c$ en unités SI.
  \item Calculer le couple électromagnétique et la puissance électromagnétique.
\end{enumerate}
\end{TLExercise}

\begin{TLExercise}{Exercice 2 — Bilan de puissance}
La MCC de l'exercice précédent possède un couple de pertes
$C_p=0{,}18\,\mathrm{N\,m}$.
\begin{enumerate}
  \item Calculer les pertes Joule d'induit.
  \item Calculer le couple utile et la puissance utile.
  \item Déterminer le rendement de la machine.
  \item Vérifier le bilan des puissances.
\end{enumerate}
\end{TLExercise}

\begin{TLExercise}{Exercice 3 — Démarrage sous tension nominale}
Une MCC possède $U_n=120\,\mathrm V$ et $R=0{,}60\,\Omega$.
Son courant nominal est $I_n=20\,\mathrm A$.
\begin{enumerate}
  \item Calculer le courant de démarrage si la tension nominale est appliquée directement.
  \item Comparer ce courant au courant nominal.
  \item Déterminer la tension maximale à appliquer au démarrage pour limiter le courant à $1{,}5I_n$.
  \item Expliquer comment un hacheur permet de réaliser cette limitation sans ajouter une forte résistance série.
\end{enumerate}
\end{TLExercise}

\begin{TLExercise}{Exercice 4 — Identification rotor bloqué}
Rotor bloqué, on applique un échelon de $12\,\mathrm V$. Le courant tend vers
$6{,}0\,\mathrm A$ et atteint $63\,\%$ de sa valeur finale après $8{,}0\,\mathrm{ms}$.
\begin{enumerate}
  \item Déterminer la résistance $R$ de l'induit.
  \item Déterminer la constante de temps électrique $\tau_e$.
  \item En déduire l'inductance $L$.
  \item Écrire l'expression de $i(t)$ après l'application de l'échelon.
\end{enumerate}
\end{TLExercise}

\begin{TLExercise}{Exercice 5 — Génératrice et freinage récupératif}
Une MCC tourne à $150\,\mathrm{rad\,s^{-1}}$ avec $K_e=0{,}32$. Elle est raccordée,
par un convertisseur réversible, à un bus continu de $42\,\mathrm V$. On néglige
l'inductance et on prend $R=0{,}30\,\Omega$.
\begin{enumerate}
  \item Calculer $E$.
  \item Déterminer le courant lorsque la machine renvoie de l'énergie vers le bus.
  \item Calculer la puissance électrique transférée au bus avec une convention de signe clairement définie.
  \item Identifier le quadrant si $\Omega>0$ et le couple électromagnétique est négatif.
\end{enumerate}
\end{TLExercise}

\begin{TLExercise}{Exercice 6 — Modèle dynamique et schéma-blocs}
On considère une MCC avec $R$, $L$, $K_e=K_c=K$, $J$ et un frottement visqueux $f$.
Le couple résistant $C_r$ est une perturbation.
\begin{enumerate}
  \item Écrire les équations électrique et mécanique dans le domaine de Laplace.
  \item Retrouver les blocs $1/(Lp+R)$ et $1/(Jp+f)$ du schéma fonctionnel.
  \item Montrer que, pour $C_r=0$, la fonction de transfert $\Omega(p)/U(p)$ vaut
  \[
    \frac{K}{(Lp+R)(Jp+f)+K^2}.
  \]
  \item Commenter l'influence de la force contre-électromotrice sur la dynamique.
\end{enumerate}
\end{TLExercise}
